\section{引言：从离散到连续}

在前几讲中，我们已经学习了扩散模型（Diffusion Models）的基础内容，包括 DDPM（Denoising Diffusion Probabilistic Models）的前向过程、反向过程，以及 DDIM（Denoising Diffusion Implicit Models）将随机生成过程转化为确定性过程以减少采样步数的方法。这些基础知识已经足够在许多应用场景中构建和训练自己的扩散模型。

本讲的核心目标是从一个全新的视角——\textbf{连续时间域}——重新理解扩散模型。我们将看到：

\begin{itemize}
    \item 离散时间的 DDPM 和 SMLD（Score Matching with Langevin Dynamics）都可以被统一到一个连续时间的 SDE（Stochastic Differential Equation，随机微分方程）框架中；
    \item 给定前向 SDE，可以直接推导出反向 SDE；
    \item 存在一个与 SDE 等价的 ODE（Ordinary Differential Equation，常微分方程）形式，称为 Probability Flow ODE；
    \item 这一连续框架为我们提供了更多设计扩散模型的自由度和理论工具。
\end{itemize}

\begin{knowledgebox}{课程背景}
本讲基于 Song et al.（2021）的经典论文 ``Score-Based Generative Modeling through Stochastic Differential Equations'' 展开。该论文统一了 SMLD（NCSN）和 DDPM 两种范式，提出了通用的 SDE 框架。这是理解现代扩散模型理论的关键文献。
\end{knowledgebox}

\begin{importantbox}{核心观点}
之前我们在\textbf{离散时间}下定义扩散模型的前向和反向过程。本讲的关键转变是：当离散时间步数 $N \to \infty$ 时，整个前向加噪过程可以用一个 SDE 来描述，而反向去噪过程同样对应一个反向 SDE。这种连续时间视角不仅统一了 SMLD 和 DDPM，还为设计新型扩散模型和高效采样器提供了理论基础。
\end{importantbox}


